(1) 令 $s=\sigma/\rho$，色散关系 $\omega=\sqrt{gk+sk^3}$。相速度 $v_p=\omega/k=\sqrt{g/k+sk}$；群速度 $v_g=d\omega/dk=(g+3sk^2)/(2\sqrt{gk+sk^3})$，两者一般不相同。
(2) 最小化相速度平方 $f(k)=g/k+sk$，导数 $f'(k)=-g/k^2+s$，故唯一驻点 $k_*=\sqrt{g/s}=\sqrt{\rho g/\sigma}$。二阶导 $f''(k)=2g/k^3>0$，因此是真正的全局最小值；$v_{p,\min}^2=f(k_*)=2\sqrt{gs}=2\sqrt{g\sigma/\rho}$。
(3) 利用一般恒等式 $\omega=kv_p(k)$，求导给 $v_g=v_p+k\,dv_p/dk$。在 $k_*$ 处相速度斜率为零，所以 $v_g(k_*)=v_p(k_*)$；这一交点由极值数学结构保证，不需要把相速度和群速度在所有波数上强行视为相等。
(4) 在 $sk^2\ll g$ 时 $\omega\simeq\sqrt{gk}$、$v_p\simeq\sqrt{g/k}$、$v_g\simeq v_p/2$；在 $sk^2\gg g$ 时 $\omega\simeq\sqrt{s}k^{3/2}$、$v_p\simeq\sqrt{sk}$、$v_g\simeq3v_p/2$。若水深有限为 $h$，频散关系应变为 $\omega^2=(gk+sk^3)\tanh(kh)$，而深水假设要求 $kh\gg1$。
